Dans le cadre de l’évaluation d’une école, la question suivante a été posée aux élèves.

 \og{} Utilisez-vous les jeux de cour ? \fg{}.
 
 \begin{itemize}
 	\item 160 élèves ont répondu à cette enquête dont 55~\% de filles. 
 	\item La moitié des filles a déclaré utiliser les jeux de cour.
 	\item Les trois quarts des garçons ont déclaré utiliser les jeux de cour.
 \end{itemize}
 
 \begin{enumerate}
 	\item Compléter le tableau suivant :
 	
 	{\renewcommand{\arraystretch}{1.8}
 	\begin{tabularx}{\linewidth}{|X|*{3}{>{\centering\arraybackslash}X|}}
	 	\hline
	 	 & Filles & Garçons & Total\\
	 	 \hline
	 	Utilisent les\newline jeux de cour & & & \\
	 	\hline
	 	N'utilisent pas\newline les jeux de cour & & & \\
	 	\hline
	 	Total & & & 160 \\
	 	\hline
 	\end{tabularx}}
	\item Calculer le pourcentage d’élèves qui utilisent les jeux de cour. 
	\item Parmi les élèves ayant déclaré utiliser les jeux de cour, quel est le pourcentage de filles ? Donner le résultat arrondi à l’unité.
 \end{enumerate}

